\documentclass[12pt,a4paper]{article}
\usepackage[utf8]{vntex, inputenc}
\usepackage{amsmath, amssymb, anyfontsize, array, color,  enumerate, enumitem, fancyhdr, float, graphicx, hyperref, lastpage, listings, mathtools, mips, natbib, pifont, tikz}
\usepackage[open,openlevel=0]{bookmark}
\usepackage[framemethod=tikz]{mdframed}

\usetikzlibrary{calc}
\newcommand\HRule{\rule{12cm}{1pt}}

\pagestyle{fancy}
\fancyhead{} % clear all header fields
\fancyhead[L]{
  \color{blue}
  \begin{tabular}{rl}
    \begin{picture}(25,15)(0,0)
    \put(0,-8){\includegraphics[width=8mm, height=8mm]{BK.png}}
    \end{picture}
    \begin{tabular}{l}
      \textbf{\bf \ttfamily Ho Chi Minh University of Technology}\\
      \textbf{\bf \ttfamily Faculty of Computer Science \& Engineering}
    \end{tabular}
  \end{tabular}
}
\fancyhead[R]{
  \begin{tabular}{l}
    \tiny \bf \\
    \tiny \bf
  \end{tabular}
}
\fancyfoot{} % clear all footer fields
\fancyfoot[L]{\scriptsize \ttfamily Computer Architecture Assignment}
\fancyfoot[R]{\scriptsize \ttfamily Page {\thepage}/\pageref{LastPage}}
\renewcommand{\headrulewidth}{0.3pt}
\renewcommand{\footrulewidth}{0.3pt}

\definecolor{dkgreen}{rgb}{0,0.6,0}
\definecolor{gray}{rgb}{0.5,0.5,0.5}
\definecolor{mauve}{rgb}{0.58,0,0.82}

\lstset{frame=tb,
  language=[MIPS]Assembler,
  aboveskip=3mm,
  belowskip=3mm,
  showstringspaces=false,
  columns=flexible,
  basicstyle={\small\ttfamily},
  numbers=none,
  numberstyle=\tiny\color{gray},
  keywordstyle=\color{blue},
  commentstyle=\color{dkgreen},
  stringstyle=\color{mauve},
  breaklines=true,
  breakatwhitespace=true,
  tabsize=3,
  numbers=left,
  stepnumber=1,
  numbersep=1pt,
  firstnumber=1,
  numberfirstline=true
}

\begin{document}
\begin{titlepage}
  \begin{tikzpicture}[remember picture, overlay]
    \draw[line width = 2pt] ($(current page.north west) + (1in,-1in)$) rectangle ($(current page.south east) + (-1in,1in)$);
  \end{tikzpicture}

  \begin{center}

    % Upper part of the page
    \textbf{\fontsize{12pt}{1pt}\selectfont HO CHI MINH CITY UNIVERSITY OF TECHNOLOGY}\\[0.5cm]
    {\fontsize{13pt}{1pt}\selectfont Faculty of Computer Science \& Engineering}\\[0.5cm]
    \begin{figure}[H]
      \centering
      \includegraphics[width=1.7in,height=1.7in]{BK.png}
    \end{figure}

    % Title
    \HRule\\[0.5cm]
    { \textbf{\fontsize{25pt}{1pt}\selectfont Computer Architecture Assignment}}\\[0.4cm]

    \HRule\\[0.8cm]
    \begin{minipage}{0.545\textwidth}
      \begin{flushleft}
        \textbf{Authors:}\\
        \begin{tabular}{l l}
          Lưu Nguyễn Hoàng Minh & 1952845 \\
          Vũ Anh Nhi            & 1952380 \\
          Nguyễn Chính Khôi     & 1952793 \\
          Nguyễn Hoàng          & 1952255 \\
        \end{tabular}
      \end{flushleft}
    \end{minipage}
    \begin{minipage}{0.4\textwidth}
      \begin{flushright}
        \textbf{Class:}\\
        CO2008\_CC02\\
        \textbf{Lecturer:}\\
        Kiều Đỗ Nguyên Bình\\

      \end{flushright}
    \end{minipage}

    \vfill

    % Bottom of the page
    \vspace{2cm}
    {\large} %{\large \today}
  \end{center}
\end{titlepage}


\pdfbookmark[0]{Question 1}{ques1}
\section*{Question 1:}
Given the following program with MIPS declaration in the data section of a MIPS program
\begin{mdframed}[hidealllines=true,backgroundcolor=magenta!10]
  \begin{lstlisting}
    .data
      nums: .word <an integer number>
      elems: .word <array elements>
  \end{lstlisting}
\end{mdframed}
Where \(<an\;integer\;number>\) will store the number of elements in the array \(elems\).
\(elems\) is an array that stores integer elements whose size is equal to value \(<an\;integer\;number>\).
You are required to choose those values when developing and testing your program.
\begin{enumerate}
  \item \textbf{Write a MIPS program that sort the array \(elems\) in ascending order using the bubble sort algorithm.}
        \begin{mdframed}[hidealllines=true,backgroundcolor=magenta!10]
          \begin{lstlisting}
    .data
      nums: .word 10
      elems: .word 1 0 2 9 3 8 4 7 5 6
      space: .asciiz " "

.text

	# load values into registers
      addi  $t9, $zero, 1       # 1
      la    $s0, elems          # Load the base address of the array arr[0] into $s0
      addi  $s1, $s0, 4         # Load the base address of the arr[1] into $s1
      add   $s2, $zero, $zero   # Index i = 0
      addi  $s3, $s2, 1         # Index j = i + 1
      la    $t7, nums           # Limiter base address
      lw    $s7, 0($t7)         # Limiter used as stopping condition in the loop for j, limiter = (n - i)

  Sort:
      beq   $t9, $s7, Exit	# If the limiter is now 1, end the program.

      lw    $t0, 0($s0)		# Load data of arr[i] into $t0
      lw    $t1, 0($s1)		# Load data of arr[j] into $t1

      slt   $t2, $t0, $t1	# Compare arr[i] and arr[j]
      beq   $t2, $t9, NoSwap	# If arr[i] < arr[j], skip to NoSwap
      sw    $t0, 0($s1)		# Else swap arr[i] and arr[j]
      sw    $t1, 0($s0)

  NoSwap:
  				# After either swap or noSwap, update i and j
      addi  $s0, $s0, 4		# Increase the address of arr[i] to arr[i+1]
      addi  $s1, $s0, 4		# Increase the address of arr[j] to arr[i+1]
      addi  $s2, $s2, 1		# Increase i = i + 1
      addi  $s3, $s2, 1		# Increase j = i + 1

      bne   $s3, $s7, NotEnd	# Continue if j < limiter
      la    $s0, elems		# Load arr[i] = arr[0] to $s0
      addi  $s1, $s0, 4		# Load arr[j] = arr[1] to $s1
      add   $s2, $zero, $zero	# Set i = 0
      addi  $s3, $s2, 1		# Set j = 1
      addi  $s7, $s7, -1	# Decrease the limiter by 1

  NotEnd:
      j Sort

  Exit:
      la    $s0, elems		# Reset all values of registers for printing
      add   $s2, $zero, $zero
      la    $t7, nums
      lw    $s7, 0($t7)

  Print:
      beq   $s2, $s7, SPrint
      lw    $s1, 0($s0)

      # Print array element
      li    $v0, 1
      move  $a0, $s1
      syscall

      # Print space
      li  $v0, 4
      la  $a0, space
      syscall

      addi  $s0, $s0, 4
      addi  $s2, $s2, 1
      j Print

  SPrint:
      # Exit
      li      $v0, 10
      syscall

          \end{lstlisting}
        \end{mdframed}


        \begin{mdframed}[hidealllines=true,backgroundcolor=magenta!10]
          \begin{lstlisting}
    Output: 0 1 2 3 4 5 6 7 8 9
         \end{lstlisting}
        \end{mdframed}
  \item \textbf{Calculate the execution time of your program if one instruction requires 1 ns for processing.} \\
        % TODO

        There are 714 instructions in total for sorting an array with 10 elemetns (including the printing instructions).

        \begin{itemize}
          \item 7: load values into registers.
          \item 576: sort the array
          \item 129: print the sorted array
          \item 2: exit instructions
        \end{itemize}

        The total execution time of our program is \(714 ns\).

\end{enumerate}

\newpage
\pdfbookmark[0]{Question 2}{ques2}
\section*{Question 2:}
\textbf{MIPS programm:}
\begin{mdframed}[hidealllines=true,backgroundcolor=magenta!10]
  \begin{lstlisting}
.data
	.align 2
	student: .space 180
	gen_select: .asciiz "Male\0|Female"
	read: .asciiz "Student "
	ID: .asciiz "Student id: "
	name: .asciiz "Student name: "
	age: .asciiz "Student age: "
	gender: .asciiz "Student gender: "
	av_score: .asciiz "Student average score: "
	msg: .asciiz "This is a list of students\n"
	sign_limit: .double 2147483647
.text

.globl main

main:
	la $s0, student

	# for(int i = 0; i != 5; i++) readStudent(i);
	addi $t5, $zero, 0
readStd:
	beq $t5, 5, readExit
	move $a2, $t5		# pass parameter std_idx
	jal readStudent
	addi $t5, $t5, 1
	j readStd
readExit:

	li $v0, 4
	la $a0, msg
	syscall

	# for(int i = 0; i != 5; i++) printStudent(i);
	addi $t5, $zero, 0
printStd:
	beq $t5, 5, printExit
	move $a2, $t5		# pass parameter std_idx
	jal printStudent
	addi $t5, $t5, 1
	j printStd
printExit:

	#sys_exit
	li $v0, 10
	syscall

readStudent:		# a2 = std_index
	addi $t7, $zero, 36 	# Students struct size
	mul $t1, $a2, $t7	# Index offset
	add $t0, $s0, $t1	# student address (address of student at index std_idx)

	#Print Student label
	li $v0, 4
	la $a0, read
	syscall
	li $v0, 1
	move $a0, $a2
	syscall
	#Print new line
	li $v0, 11
	li $a0, 0xA
	syscall

	#Read ID
	li $v0, 4
	la $a0, ID
	syscall
	li $v0, 7
	syscall
	#if exceeds sign limit
	l.d $f2, sign_limit
	addi $t3, $zero, 0
	c.lt.d $f0, $f2
	bc1t no_exceed
	sub.d $f2, $f0, $f2
	cvt.w.d $f2, $f2
	mfc1 $t3, $f2
no_exceed:
	cvt.w.d $f0, $f0
	mfc1 $t2, $f0
	addu $t2, $t2, $t3
	#Store ID - 1 Word (4 Bytes)
	sw $t2, 0($t0) 		# ID address = student

	#Read & store name - 25 Bytes
	la $a0, name
	li $v0, 4
	syscall
	addi $a0, $t0, 4	# name address = ID address + sizeof(ID) (a0 = address buffer)
	li $a1, 25		# no more than 24 characters with 25th character being null-terminator
	li $v0, 8
	syscall

	#Read age
	la $a0, age
	li $v0, 4
	syscall
	li $v0, 5
	syscall
	#Store age - 1 Byte
	addi $t1, $v0, 0
	sb $t1, 29($t0)		# age address = name address + sizeof(name)

	#Read gender
	la $a0, gender
	li $v0, 4
	syscall
	li $v0, 5
	syscall
	#Store gender - 1 Byte
	addi $t1, $v0, 0
	sb $t1, 30($t0)		#gender address = age address + sizeof(age)

	#Read average score
	la $a0, av_score
	li $v0, 4
	syscall
	li $v0, 6
	syscall
	#Store average score - 1 Word (4 Bytes)
	swc1 $f0, 32($t0)	#av_score address = gender address + sizeof(gender) + 1 Padding Byte

	# Total Bytes used: 35
	# Padding 1 Bytes, total Bytes occupied: 36

	#Print new line
	li $v0, 11
	li $a0, 0xA
	syscall

	jr $ra

printStudent: 		# a2 = std_index
	addi $t7, $zero, 36 	# Students struct size
	mul $t1, $a2, $t7	# Index offset
	add $t0, $s0, $t1	# student address (address of student at index std_idx)

	#Load ID - 1 Word (4 Bytes)
	lw $t2, 0($t0) 		# ID address = student
	#Print ID
	li $v0, 4
	la $a0, ID
	syscall
	li $v0, 36		# Print unsigned int
	move $a0, $t2
	syscall

	#Print new line
	li $v0, 11
	li $a0, 0xA
	syscall

	#Load & print name - 25 Bytes
	la $a0, name
	li $v0, 4
	syscall
	li $v0, 4
	addi $a0, $t0, 4	# name address = ID address + sizeof(ID) (a0 = address buffer)
	syscall

	#Load age - 1 Byte
	lb $t1, 29($t0)		# age address = name address + sizeof(name)
	#Print age
	la $a0, age
	li $v0, 4
	syscall
	li $v0, 1
	move $a0, $t1
	syscall

	#Print new line
	li $v0, 11
	li $a0, 0xA
	syscall

	#Load gender - 1 Byte
	lb $t1, 30($t0)		#gender address = age address + sizeof(age)
	#gen_select handling
	la $t2, gen_select
	beq $t1, 0, jump
	addi $t1, $zero, 6
jump:
	add $t1, $t1, $t2
	#Print gender
	la $a0, gender
	li $v0, 4
	syscall
	li $v0, 4
	move $a0, $t1
	syscall

	#Print new line
	li $v0, 11
	li $a0, 0xA
	syscall

	#Load average score - 1 Word (4 Bytes)
	lwc1 $f0, 32($t0)	#av_score address = gender address + sizeof(gender)
	#Print average score
	la $a0, av_score
	li $v0, 4
	syscall
	li $v0, 2
	mov.s $f12, $f0
	syscall

	#Print new line
	li $v0, 11
	li $a0, 0xA
	syscall

	#Print new line
	li $v0, 11
	li $a0, 0xA
	syscall

	jr $ra

      \end{lstlisting}
\end{mdframed}
\begin{enumerate}
  \newpage
  \item \textbf{Organize memory allocation for the \textbf{Students struct} (padding is required).} \\

        Struct Students allocates total of 36 bytes of memory, its member:
        \begin{itemize}
          \item id: occupies from byte index 0 to 3. Total of 4 bytes.
          \item name[25]: occupies from byte index 4 to 28, with maximum 24 characters and one null-terminator at the end. Total of 25 bytes.
          \item age: occupies byte index 29. Total of 1 byte.
          \item is\_male: occupies byte index 30. Total of 1 byte.
          \item Padding 1 bytes index 31 for storing next member with size of one word.
          \item average\_score: occupies from byte index 32 to 35. Total of 4 bytes.
        \end{itemize}

        With the size of array students is five, we allocate total of \(36 \cdot 5 = 180\) bytes of memory

  \item \textbf{Initialize the array of 5 students. Assign any value for their information on the main program.} \\

        Get students information through user input:\\
        Student 0\\
        Student id: 2500000000 \\
        Student name: Peter \\
        Student age: 20\\
        Student gender: 1\\
        Student average score: 7.81\\
        \\
        Student 1\\
        Student id: 64652156\\
        Student name: Daisy \\
        Student age: 19\\
        Student gender: 0\\
        Student average score: 9.75\\
        \\
        Student 2\\
        Student id: 3548896568\\
        Student name: Charles\\
        Student age: 26\\
        Student gender: 1\\
        Student average score: 8.34\\
        \\
        Student 3\\
        Student id: 527846889\\
        Student name: Katherine\\
        Student age: 23\\
        Student gender: 0\\
        Student average score: 6.23\\
        \\
        Student 4\\
        Student id: 1565486232\\
        Student name: Dick\\
        Student age: 24\\
        Student gender: 1\\
        Student average score: 5.16\\

  \item \textbf{Write a MIPS program for print\_student procedure.}\\

        Included in the MIPS program above with label printStudent.


  \item \textbf{Use print\_student procedure to print information of 5 assigned students.}\\

  \item[\(\bullet\)] Result:\\

        \begin{mdframed}[hidealllines=true,backgroundcolor=magenta!10]
          \begin{lstlisting}

    This is a list of students
    Student id: 2500000000
    Student name: Peter
    Student age: 20
    Student gender: 1
    Student average score: 7.81

    Student id: 64652156
    Student name: Daisy
    Student age: 19
    Student gender: 0
    Student average score: 9.75

    Student id: 3548896568
    Student name: Charles
    Student age: 26
    Student gender: 1
    Student average score: 8.34

    Student id: 527846889
    Student name: Katherine
    Student age: 23
    Student gender: 0
    Student average score: 6.23

    Student id: 1565486232
    Student name: Dick
    Student age: 24
    Student gender: 1
    Student average score: 5.16

  \end{lstlisting}
        \end{mdframed}
\end{enumerate}

\end{document}
